\chapter{Case Selection and Introduction}
The case selected is the 2020 Twitter Bitcoin Scam and the incident occurred on July 15, 2020, where 130 high-profile Twitter accounts were compromised by unauthorized parties to promote a Bitcoin scam. The accounts compromised included famous people such as former US presidents Barack Obama and Joe Biden, Billionaires Bill Gates and Elon Musk as well as other internationally known people \cite{twitterhack, isacablog}.
\newline
The attack can be classified as a social engineering attack, where attackers trick victims by exploiting human emotions such as fear, trust or healpfulness to help them gain access to systems \cite{erbschloe2004trojans}. 
\newline
In the context of this incident, the attackers tricked Twitter employees with administrative rights to give them their credentials and used these credentials to access Twitter's internal systems.
\newline 
By July 30, 2020, Twitter confirmed that the attackers first tricked employees with lower rights into giving them their credentials and used the these employee accounts to get the credentials from employee accounts that were authorized to use the administrative tools.
\newline
This incident is significant because it showed that the nature of cyberattacks are not just technical anymore, but also an attack on human behaviour as the attackers managed to engage in this attack purely by social engineering without the need of exploiting technical vulnerabilities. 
\newline
Another thing that needs to be highlighted is the rapid pace as the incident just lasted for about two hours before the compromised accounts were restored and yet, the attackers still managed to caused considerable damage.
\newline
I chose this incident not only because of the two points mentioned before why the incident is significant, but also because I have a personal interest in cryptocurrency such as Bitcoin. In a way, I am fascinated in how one benefit in cryptocurrency, anonymity, is used by attackers to hide their own identities \cite{kethineni2020rise}.